\documentclass{article}
\usepackage{amsmath}
\usepackage{graphicx} % Required for inserting images
\usepackage{xcolor}
\usepackage{algorithmic}
\usepackage{algorithm}
% \usepackage{algpseudocode}
\usepackage{booktabs}
\title{Inventory Routing}
\date{March 2024}

\begin{document}

\maketitle
% \section{TO DO}
% \begin{itemize}
%     \item Add a selection mechanism to the local search rather than performing a full enumeration of the neighbourhood.
%     \item Add a disturbance mechanism if required.
%     \item Define parameters used in the algorithm and tune with iRace
%     \item In parallel, documenting the implementation.
% \end{itemize}

\section{Introduction}


\section{Literature review}


\section{Problem description}


\section{Description of the Algorithm}


\begin{algorithm}[H]
\caption{\textbf{Main algorithm}}
\begin{algorithmic}[1]
\STATE $i\gets 1$
\STATE $s^* \gets$ Empty \AND  $f(s^*) \gets +\infty$ 
\WHILE{$i\leq 5$} 
\STATE $s\gets$ Construct initial schedule 
\STATE $s\gets$ Hybrid\_Genetic\_Search$(s)$ \COMMENT{Improve initial routing} 
\IF{$f(s)<f(s^*)$}
    \STATE $s^*\gets s$
\ENDIF
\STATE Disturbance\_Counter $\gets 1$
\WHILE{Disturbance\_Counter $\leq 5$}
    \STATE $s\gets$ Local\_Search$(s)$
    \STATE Disturbance\_Counter $\gets$ Disturbance\_Counter$+1$
    \IF{$s$ is feasible}
        \STATE $s\gets$ Rebalance$(s)$
        \IF{criterion meet}
            \STATE Half the penalty
        \ENDIF
        \IF{$(s)<f(s^*)$}
            \STATE $s^*\gets s$
        \ENDIF
    \ELSIF{criterion meet}
        \STATE Double the penalty
    \ENDIF
    \STATE $s \gets$ Disturbance$(s^*,$ Disturbance\_Counter$)$  
\ENDWHILE
\STATE $s^*\gets$ Hybrid\_Genetic\_Search$(s^*)$ %Skip if s^* is not updated

%Giving one vehicle instance, disturbance removing schedule for the whole day. Could extend it to instances with multiple vehicles.
%
\STATE $i\gets i+1$
\ENDWHILE
\RETURN solution $s^*$
\end{algorithmic}
\end{algorithm}

\textcolor{red}{CA: THE FOLLOWING IS A SIMPLER PSEUDO-CODE I USED IN THE PRESENTATION}

% \begin{algorithm}[H]
% \caption{\textbf{ILS-LR}}
% \begin{algorithmic}[1]
% %\STATE $i\gets 1$
% \STATE $s^* \gets$ Empty \AND  $f(s^*) \gets +\infty$ 
% \WHILE{$maxIter$} 
% \STATE $s\gets$ Construct initial schedule 
% \STATE $s\gets$ HGS$(s)$ \COMMENT{Improve initial routing} 
% \STATE $s^*\gets s$
% %\STATE Local Search Counter $\gets 1$
% \WHILE{$\leq maxLS$}
%     \STATE $s\gets$ Local\_Search$(s)$
%     \IF{$s$ is feasible}
%         \STATE $s\gets$ Rebalance$(s)$
%         \IF{$f(s)<f(s^*)$}
%             \STATE $s^*\gets s$
%         \ENDIF
%     \ENDIF
%     \IF{criterion met}
%         \STATE Update Penalty
%     \ENDIF
% \ENDWHILE
% \STATE $s^*\gets$ HGS$(s^*)$ %Skip if s^* is not updated

% %Giving one vehicle instance, disturbance removing schedule for the whole day. Could extend it to instances with multiple vehicles.
% %
% \ENDWHILE
% \RETURN solution $s^*$
% \end{algorithmic}
% \end{algorithm}





\subsection{Initial feasible solution}




% The goal of this procedure is to construct a feasible solution. Although we expect this procedure does not produce a solution with a good logistic ratio, this procedure can help the heuristic to set the initial upper bound. 

The algorithm described in this section constructs initial feasible solutions. These solutions are used as starting solutions for a local search algorithm.

Minimum total delivery quantity for customer $i\in N'$ $= \sum^{T}_{t=1}d_{i,t}-I_{i,0}$

Approximate number of visits $v_i= \frac{\sum^T_{t=1}d_{i,t}-I_{i,0}}{U_i}$, $v_i= \frac{\sum^T_{t=1}d_{i,t}-I_{i,0}}{\min{I_{i,t}|}}$

Minimum delivery quantity per visit $=\frac{\sum^T_{t=1}d_{i,t}-I_{i,0}}{v_i}$





Let $C_t$ be a set of customers who have stock-out at period $t\in T$. Let $R$ be a set of non-overlapping grids. Each grid is defined by four numbers, $Minimum_xCoord + \frac{Maximum_xCoord-MinimumxCoord}{D_v}$, $Minimum_yCoord +\frac{Maximum_yCoord-Minimum_yCoord}{D_h}$.


Let $B_{t,r}$ be sets of customers who have first stock-out at time $t$ and have the location within grid $r$. For each $t\in T$, $B_{t,r}$ are sorted by the distance between the centre point of grid $r$ to the depot. For each $t\in T$ and $r\in R$, customers in $B_{t,r}$ is sorted by a decreasing order with respect to the demand.



\begin{algorithm}[H]
\caption{\textbf{Construct initial schedule}}
\begin{algorithmic}[1]
\STATE Initialize $s$ to empty solution
\REPEAT
\STATE Randomly choose customer $i$ in $B_{t,r}$
\STATE $t_{current}\gets 1$ 
\STATE $t_{stock\_out} \gets$ Next stock out period
\REPEAT
\STATE $s\gets $Assign a visit to customer $i$ at a random time period $t$ between $t_{current}$ and $t_{stock\_out}$
\STATE $s\gets $ Assign a vehicle to this visit with the maximum possible delivery quantity
\STATE Update inventory level  
\STATE $t_{stock\_out}\gets$ Next stock out period
\STATE $t_{current}\gets t$
\IF{ $t_{current}$ = $t_{stock\_out}$ \AND stock out remains}
    \IF{Customer $i$ is visited at every time period up to $t_{stock\_out}$ \AND stock out remains}
        \STATE Restart the procedure from line 1
    \ELSE
        \STATE $t_{current}\gets$ Time period before $t$ with no visit to customer $i$ 
        \STATE $s\gets$ Assign a visit at $t_{current}$
        \STATE $s\gets $ Assign a vehicle to this visit with the maximum possible delivery quantity
        \STATE Update inventory level  
        \STATE $t_{stock\_out}\gets$ Next stock out period
    \ENDIF
\ENDIF
\UNTIL{No stock out for customer $i$}
\UNTIL{No stock out for all customers}
\RETURN solution 
\end{algorithmic}
\end{algorithm}

% \begin{algorithm}[H]
% \caption{\textbf{Construct initial schedule}}
% \begin{algorithmic}[1]
% \STATE Initialize $s$ to empty solution
% \REPEAT
% \STATE Randomly choose customer $i$ in $B_{t,r}$
% \STATE $t_{current}\gets 1$ 
% \STATE $t_{stock\_out} \gets$ Next stock out period
% \REPEAT
% \STATE $s\gets $Assign a visit to customer $i$ at a random time period $t$ between $t_{current}$ and $t_{stock\_out}$
% \STATE $s\gets $ Assign a vehicle to this visit with the maximum possible delivery quantity
% \STATE Update inventory level  
% \STATE $t_{stock\_out}\gets$ Next stock out period
% \STATE $t_{current}\gets t$
% \IF{ $t_{current}$ = $t_{stock\_out}$ \AND stock out remains}
%     \IF{Customer $i$ is visited at every time period up to $t_{stock\_out}$ \AND stock out remains}
%         \STATE Restart the procedure from line 1
%     \ELSE
%         \STATE $t_{current}\gets$ Time period before $t$ with no visit to customer $i$ 
%         \STATE $s\gets$ Assign a visit at $t_{current}$
%         \STATE $s\gets $ Assign a vehicle to this visit with the maximum possible delivery quantity
%         \STATE Update inventory level  
%         \STATE $t_{stock\_out}\gets$ Next stock out period
%     \ENDIF
% \ENDIF
% \UNTIL{No stock out for customer $i$}
% \UNTIL{No stock out for all customers}
% \RETURN solution $s$
% \end{algorithmic}
% \end{algorithm}


%\subsection{Ideas to test}
%\begin{itemize}
    %\item Grouping the customers into different groups based on the location to the depot. Let set $G_{x,y,s}$ be a set of groups. Index $x$ indicate how many groups are splited based the difference on the x-coordinate between the customer and the depot, similarly, index $y$ indicate how many groups are splited based the difference on the y-coordinate between the customer and the depot. Index $s$ indicate the number of groups splited via the stock out period. Taking the minimum and maximum difference of x-coordinate or y-coordinate, based on this range and split the customers into some number of groups. In each group, customers are partitioned based the first stock out period. Then the initial procedure randomly pick one customer from the first partition and repeat until all customers are considered.
%\end{itemize}

\subsection{Local search}

The local search uses three neighbourhood operators, denoted by $O_1$, $O_2$, and $O_3$. 

\begin{itemize}
    \item Neighbourhood operator $O_1$ swaps or inserts customers between two routes on the same day, followed by an update of the corresponding delivery quantities and inventory levels.
    
    \item Neighbourhood operator $O_2$, also applied on a single day. It swaps customers in a route with unallocated customers, removes them from a route, or inserts unallocated customers into a route. As with $O_1$, the delivery quantities and inventory levels are updated accordingly.
    
    \item Neighbourhood operator $O_3$ focuses on a single vehicle route by relocating customers to different positions within the same route. Please note that this operator does not change the delivery quantity and inventory level. 
\end{itemize}

The local search is an iterative procedure starting from the input solution $s$. It repeatedly and exhaustively explores solutions defined by the neighbourhood operators $O_1$, $O_2$, and $O_3$ using neighbourhood functions $N_1$, $N_2$, and $N_3$. The search terminates when no operator can further improve the solution.  

\begin{algorithm}[H]
\caption{\textbf{Local search$(s)$}}
\begin{algorithmic}[1]
\STATE True\_local $\gets$ \FALSE
\WHILE{True\_local is \FALSE}
    \STATE True\_local $\gets$ \TRUE
    \FOR{$i\gets 1; i\leq 3; i\gets i+1$}
        \REPEAT
            \STATE $\hat{s}\gets s$
            \STATE $s\gets N_i(s)$
            \IF{$f(\hat{s}) \neq f(s)$}
                \STATE True\_local $\gets$ \FALSE
            \ENDIF
        \UNTIL{$f(\hat{s}) == f(s)$}
    \ENDFOR
\ENDWHILE
\end{algorithmic}
\end{algorithm}




\subsection{Rebalance delivery quantity}

This operator rebalances delivery quantities across routes on all days to optimise the logistic ratio while avoiding stockouts.



In summary, customers prioritised for additional delivery are those with:
\begin{itemize}
    \item Low current inventory (approaching danger zone)
    \item High daily demand
    \item Long gaps until next visit
    \item Infrequent visit patterns
\end{itemize}

\subsection{Perturbation}

\begin{itemize}
    \item  Randomly remove or some criterias to completely remove a route, and the quantity is removed.
\end{itemize}


\begin{itemize}
    \item  Step 1: Randomly remove a vehicle where all vehicles across all days have an equal chance. 
    \item Randomly remove or some criterias to completely remove a route, and the quantity is removed.
\end{itemize}

\subsection{Minimum and Maximum number of visits}
Before running the heuristic, we compute a customer-wise range for the number of visits using a fast inventory-based preprocessing step. Since the heuristic adopts a fill-up-to-maximum policy, the maximum quantity deliverable to customer \(i\) in a single visit is \(C_i=\min\{U_i,Q\}\). A lower bound on the number of visits required by customer \(i\) is obtained from cumulative demand coverage:
\[
LB_i=\max_{t\in T}\left\lceil\frac{\left(\sum_{\tau=1}^{t} r_{i\tau}-I_{i0}\right)^+}{C_i}\right\rceil,
\]
where $(x)^+=max(0,x)$.  

An upper bound is obtained by simulating a policy that visits the customer whenever positive delivery is possible under the fill-up rule. These bounds are computed to restrict the visit frequencies explored by the heuristic.


\begin{algorithm}[H]
\caption{\textbf{Compute customer-wise visit bounds}}
\label{alg:visit_bounds}
\begin{algorithmic}[1]
\STATE Initialize $(LB_i,UB_i)$ for all $i\in N'$
\FOR{each customer $i\in N'$}
    \STATE $C_i \gets \min\{U_i,Q\}$
    
    \STATE \textbf{// Lower bound computation//}
    \STATE $LB_i \gets 0$
    \STATE $D \gets 0$
    \FOR{$t=1$ to $H$}
        \STATE $D \gets D + r_{it}$
        \STATE $LB_i \gets \max\left\{LB_i,\left\lceil \dfrac{\max\{0,D-I_{i0}\}}{C_i}\right\rceil\right\}$
    \ENDFOR
    
    \STATE \textbf{// Upper bound computation//}
    \STATE $UB_i \gets 0$
    \STATE $I \gets I_{i0}$
    \FOR{$t=1$ to $H$}
        \IF{$I < U_i$}
            \STATE $\delta \gets \min\{C_i,U_i-I\}$
            \IF{$\delta > 0$}
                \STATE $I \gets I + \delta$
                \STATE $UB_i \gets UB_i + 1$
            \ENDIF
        \ENDIF
        \STATE $I \gets \max\{I - r_{it},0\}$
    \ENDFOR
\ENDFOR
\RETURN $(LB_i,UB_i)$  for all $i\in N'$
\end{algorithmic}
\end{algorithm}

\section{Computational results}

\begin{itemize}
    \item Add results on small benchmark instances.
    \item Limiting algorithm to 60 seconds. Also experiment with limiting the time limit for HGS
    \item Test performance without HGS. or remove some of the mechanisms.

    \item Compare explicitly with travel distance and delivery quantity

\end{itemize}





\appendix

\section{Mathematical model}\label{math model}

Let $G = (N, E)$ be a complete undirected graph where $N = \{0\} \cup N'$ is the set of locations (nodes) and $E$ is the set of edges between locations. Node $0$ represents the depot (supplier location) and $N'$ denotes the set of customers. A cost $c_{ij}$ is associated with each edge $\langle i, j \rangle$, satisfying the triangle inequality. The planning horizon $T = \{1, \dots, H\}$ consists of $H$ periods (days). 

Each customer $i \in N'$ has an inventory capacity $U_i$ and a daily demand $r_{it}$ for each period $t \in T$. The quantity produced at the depot during period $t$ is $r_{0t}$, and $I_{i0}$ denotes the initial inventory level at location $i \in N$. A fleet $K$ of homogeneous vehicles, each of capacity $Q$, is available to distribute goods from the depot to the customers.

The Inventory Routing Problem with Logistic Ratio (IRP-LR) aims to determine a distribution plan that minimizes the logistic ratio while satisfying all customer demands and the following constraints:
\begin{enumerate}
    \item Vehicle capacity is not exceeded.
    \item No stockout occurs at any node in any period.
    \item The maximum inventory level at customers is not exceeded.
\end{enumerate}

The arc-flow formulation for the IRP-LR from \cite{archetti2019exact} requires the following variables:
\begin{itemize}
    \item $I_{it}$: Inventory level at location $i$ at the end of period $t$;
    \item $y_{ij}^{kt}$: Number of times edge $\langle i, j \rangle$ is traversed by vehicle $k$ in period $t$;
    \item $q_{it}^k$: Quantity delivered to customer $i$ by vehicle $k$ in period $t$;
    \item $z_{it}^k$: Binary variable equal to 1 if customer $i$ is visited by vehicle $k$ in period $t$.
\end{itemize}

For a subset of customers $S \subseteq N'$, let $E(S) = \{\langle i, j \rangle \in E \mid i, j \in S\}$. Then, the IRP-LR can be formulated as follows:

\begin{align}
    \min \quad & \frac{\displaystyle \sum_{k \in K} \sum_{\langle i,j\rangle \in E} \sum_{t \in T} c_{ij} y_{ij}^{kt}}
    {\displaystyle \sum_{k \in K} \sum_{i \in N'} \sum_{t \in T} q_{it}^k} \label{obj} \\
    \text{s.t.} \quad 
    & I_{0t} = I_{0,t-1} + r_{0t} - \sum_{k \in K} \sum_{i \in N'} q_{it}^k, && t \in T, \label{supbal} \\
    & I_{it} = I_{i,t-1} - r_{it} + \sum_{k \in K} q_{it}^k, && i \in N', \; t \in T, \label{custbal} \\
    & I_{it} \ge 0, && i \in N, \; t \in T, \label{no_stockout} \\
    & \sum_{k \in K} q_{it}^k \le U_i - I_{i,t-1}, && i \in N', \; t \in T, \label{maxinv} \\
    & q_{it}^k \le U_i z_{it}^k, && i \in N', \; k \in K, \; t \in T, \label{link_qz} \\
    & \sum_{i \in N'} q_{it}^k \le Q z_{0t}^k, && k \in K, \; t \in T, \label{vehcap} \\
    & \sum_{k \in K} z_{it}^k \le 1, && i \in N', \; t \in T, \label{nosplit} \\
    & \sum_{j: \langle i,j \rangle \in E} y_{ij}^{kt} = 2 z_{it}^k, && i \in N, \; k \in K, \; t \in T, \label{routing1} \\
    & \sum_{\langle i,j \rangle \in E(S)} y_{ij}^{kt} \le |S| - 1, && S \subset N', \; s \in S, \; k \in K, \; t \in T, \label{subtour} \\
    & z_{it}^k \in \{0,1\}, && i \in N, \; k \in K, \; t \in T, \label{binary_z} \\
    & q_{it}^k \ge 0, && i \in N, \; k \in K, \; t \in T, \label{nonneg_q} \\
    & y_{ij}^{kt} \in \{0,1,2\}, && \langle i,j\rangle \in E, \; k \in K, \; t \in T. \label{domain_y}
\end{align}

The objective function~\eqref{obj} minimizes the \emph{logistic ratio}, defined as the ratio between total routing cost and total quantity delivered.  
Constraints~\eqref{supbal}–\eqref{custbal} represent inventory balance equations for the supplier and customers.  
Constraint~\eqref{no_stockout} ensures that no stockout occurs at any node.  
Constraint~\eqref{maxinv} enforces customer inventory capacity limits.  
Constraint~\eqref{link_qz} links delivery quantities and visit decisions.  
Constraint~\eqref{vehcap} ensures that vehicle capacities are not exceeded.  
Constraint~\eqref{nosplit} prevents split deliveries.  
Constraints~\eqref{routing1}–\eqref{subtour} ensure valid vehicle routing, with subtour elimination enforced dynamically due to its exponential number of constraints.  
Finally, \eqref{binary_z}–\eqref{domain_y} define variable domains.


\begin{align}
    \min \quad & \frac{\displaystyle \sum_{k \in K} \sum_{\langle i,j\rangle \in E} \sum_{t \in T} c_{ij} y_{ij}^{kt}}
    {\displaystyle \sum_{k \in K} \sum_{i \in N'} \sum_{t \in T} q_{it}^k} \label{obj} \\
    \text{s.t.} \quad 
    & I_{0t} = I_{0,t-1} + r_{0t} - \sum_{k \in K} \sum_{i \in N'} q_{it}^k, && t \in T, \label{supbal} \\
    & I_{it} = I_{i,t-1} - r_{it} + \sum_{k \in K} q_{it}^k, && i \in N', \; t \in T, \label{custbal} \\
    & I_{it} + \mu_{it} \ge 0, && i \in N, \; t \in T, \label{no_stockout} \\    
    &\sum_{i\in N'}\sum_{t\in T} \mu_{it} \leq 0 \label{zero stockout} \\    
    & \mu_{it} \ge 0, && i \in N, \; t \in T, \label{nonneg_q} \\
    & \sum_{k \in K} q_{it}^k \le U_i - I_{i,t-1}, && i \in N', \; t \in T, \label{maxinv} \\
    & q_{it}^k \le U_i z_{it}^k, && i \in N', \; k \in K, \; t \in T, \label{link_qz} \\
    & \sum_{i \in N'} q_{it}^k \le Q z_{0t}^k, && k \in K, \; t \in T, \label{vehcap} \\
    & \sum_{k \in K} z_{it}^k \le 1, && i \in N', \; t \in T, \label{nosplit} \\
    & \sum_{j: \langle i,j \rangle \in E} y_{ij}^{kt} = 2 z_{it}^k, && i \in N, \; k \in K, \; t \in T, \label{routing1} \\
    & \sum_{\langle i,j \rangle \in E(S)} y_{ij}^{kt} \le |S| - 1, && S \subset N', \; s \in S, \; k \in K, \; t \in T, \label{subtour} \\
    & z_{it}^k \in \{0,1\}, && i \in N, \; k \in K, \; t \in T, \label{binary_z} \\
    & q_{it}^k \ge 0, && i \in N, \; k \in K, \; t \in T, \label{nonneg_q} \\
    & y_{ij}^{kt} \in \{0,1,2\}, && \langle i,j\rangle \in E, \; k \in K, \; t \in T. \label{domain_y}
\end{align}

\bibliographystyle{elsarticle-harv} 
\bibliography{reference}
\end{document}
 # Different stopping criterions and see which triggers more frequent.
 # Observe the trajectory of the solutions
 # Compare small instances with optimal solution
 # 